\paragraph{The two traces of the mixed term.} Put \(D=H_*D_*-S_*\). Cyclicity of the trace gives the
following identity at the parent. \[
 Q_c=\tfrac12\sum_{u,v}(S_*)_{uv}(m^2D)_{vu}
       +\tfrac1{4m^2}\sum_{u,v}(m^2D)_{uv}(m^2D)_{vu}.
\] Each square in \(F\setminus\{c\}\) has three retained lines, giving six ordered pairs in the first
sum. The second sum integrates over two free lines. To express \(m^2D\) in terms of the profile, write
\(\mathcal K(v,w)=E(v,w)+\Pi_{vw}\), \(\mathcal V(v,j)=\pesodirecto(v\cap j)+\Pi_{vj}\), \(j\in
\familias(c)\). For retained lines \(v,u\), set \(\mathcal Y(v,u)=-\pesodirecto(z)T_z\) if they intersect
at \(z\in E'\), and set it to zero otherwise. Here \(T_z\) uses the same indicators as
\(\pesodirecto(z)\). Finally, for retained \(w\ne u\), set \(\mathcal S(w,u)=-\pesodirecto(i')\) when
\(w\cap u=i'\in F\setminus\{c\}\) and zero on all other pairs. Define \begin{align*}
 \mathcal B(v,u)&=\mathcal Y(v,u)+\Pi_{vu}
                    +\sum_{j\in \familias(c)}\mathcal V(v,j)\mathcal V(u,j),\\
 \mathcal E_2(v,u)&=\mathcal B(v,u)+\sum_w\mathcal K(v,w)\mathcal S(w,u)
                     +\int\mathcal K(v,w)\mathcal B(w,u)\,d\rho(w).
\end{align*} With \(B_*=D_*-S_*\) and \(K_*=H_*-I\), these three terms come from
\(m^2D=m^2B_*+(mK_*)(mS_*)+m^{-1}(mK_*)(m^2B_*)\). The discrete sum runs over each \(i'=u\cap j\), \(j\in
\familias(c)\), and each of its two transverse lines \(w\). Its availability is included in
\(\pesodirecto(i')\). In the integral we mark \(w\), since its availability is already included in
\(d\rho(w)\).

Write \(R(i)=\{(u,v):u,v\in \familias(i)\setminus\{\lambda\},\ u\ne v\}\) for a square on fiber
\(\lambda\). Applying \(\psi\) to the monomials gives \begin{align*}
 L_1&=-\tfrac12\left\langle\int_F
          \sum_{(u,v)\in R(i)}\mathbb E\mathcal E_2(v,u)\right\rangle_c,\\
 L_2^{\rm int}&=\tfrac14\left\langle\iint
 \mathbb E\big[\mathcal E_2(u,v)\mathcal E_2(v,u)-\mathcal Y(u,v)\mathcal Y(v,u)\big]
                 \,d\rho(u)\,d\rho(v)\right\rangle_c .
\end{align*} The two copies of an internally integrated line have independent free variables. We subtract
the \(\mathcal Y\mathcal Y\) block before integrating. On the board, each \(i\in E'\) contributes twelve
ordered pairs to the \(YY\) block, giving \(+3m^2\sum_{E'}y_i^2\), which cancels the corresponding term
in \(P_c\) exactly. The same cancellation can be checked in the continuum. For each pair of families,
weighting the measure of valid intersections by their availability gives \(U_DU_A\,di\), including the
sheets and the Jacobian of the intersection. Multiplying by \(W(i)^2\) gives \(\nu_1(i)W(i)\).

\paragraph{Extracting \(\Delta_1\) and \(\Delta_2\).} The tabulated rows are \(Q_1=L_1-\Delta_1\) and
\(Q_2=L_2^{\rm int}-\Delta_2\). We give the corrections separately so that each entry can be
reconstructed. In the sums below, an intersection must lie in the square and belong to the indicated
parity class. Two lines must agree identically in the free parameters of that class to count as equal.
Accidental coincidences form sets of measure zero.

For \(\Delta_1\), fix \(c,i,\lambda,(u,v)\in R(i)\) as in \(L_1\) and integrate a free line \(w\) against
\(d\rho(w)\). Set \(z=v\cap w,\ p=w\cap u\). For each \(a\in \familias(p)\setminus\{w,u\}\) and \(j\in
\familias(c)\), set \(t=a\cap j\). Mark \(\familias(c)\cup \familias(i)\cup\{w\}\). The direct monomial
in \(\mathcal K(v,w)\mathcal Y(w,u)\) contributes the difference \[
 \Delta_1=\tfrac12\left\langle\int_F\sum_{(u,v)\in R(i)}
 \int\sum_{a,j} W(z)W(p)W(t)
 \big[\psi(\{z,p,t\})-\psi(\{z\})\psi(\{p,t\})\big]\,d\rho(w)\right\rangle_c .
\] The three intersections are distinct, with \(p\in E'\). The difference counts the common transverse
line of \(z,t\) only once. If this line is \(\ell\), the bracket equals
\(\psi(\{z\})\psi(\{p,t\})(U_\ell^{-1}-1)\). The three direct parts give a negative sign in the product
in the first trace before multiplication by its prefactor \(-1/2\). For example, in relative coordinates,
taking \(c=(0,0)\), \(i=(0,s)\), \(u=C_s\), \(v=D_{-s}\), \(w=R_h\), \(a=A_{h+s}\), \(j=R_0\) gives
\(z=(h,h+s)\), \(p=(h,s)\), \(t=(0,h+s)\). Column \(h+s\) occurs in both \(z\) and \(t\).

For \(\Delta_2\), mark \(\familias(c)\cup\{u,v\}\), with \(u,v\) free. When the intersection exists, the
weight of its transverse lines is \[
 K_{ab}=W(a\cap b)\prod_{\ell\in \familias(a\cap b)\setminus\{a,b\}}U_\ell .
\] For each \(j\in \familias(c)\), take \(p_2=v\cap j\), \(w_2\in \familias(p_2)\setminus\{v,j\}\),
\(t_2=u\cap w_2\), \(k_2=K_{vj}K_{uw_2}\). The two possible patterns for the other two vertices give \[
\begin{array}{c|c|c|c}
 &p_1,\ t_1 &k_1&\text{multiplicity }s\\ \hline
 \mathcal K\mathcal S\ \text{times}\ \mathcal K\mathcal S&
 u\cap j,\ v\cap w_1 & K_{uj}K_{vw_1}&1\\
 \mathcal Y\ \text{times}\ \mathcal K\mathcal S&
 v\cap u,\ a\cap j & K_{vu}K_{aj}&2
\end{array}
\] In the first row, \(w_1\) ranges over \(\familias(p_1)\setminus\{u,j\}\). In the second, \(a\) ranges
over \(\familias(p_1)\setminus\{v,u\}\). Each pattern retains four distinct vertices. For each pair of
intersections defining \(k_d\), \(d=1,2\), collect the transverse lines of its two intersections, namely
the lines whose factors \(U\) occur in \(k_d\). If the two sets share an unmarked line \(\ell\), include
the following summand. The four intersections then form a square and the transverse line is unique. \[
 \Delta_2=\tfrac14\left\langle\iint\sum_{\text{admissible patterns}}
          s\,k_1k_2(U_\ell^{-1}-1)\,d\rho(u)\,d\rho(v)\right\rangle_c .
\] The list follows by expanding the products in \(\mathcal E_2\). A weight \(\pesodirecto\) contributes
the four coordinates of its square after the marked lines are excluded. To compare two labels of the same
family, subtract their affine expressions in the free coordinates. Only an identically zero difference
contributes the factor \(U_\ell^{-1}-1\). It suffices to compare the constant term and the coefficients
of each parameter in the four labels \(r,q,r-q,r+q\). For \(\Delta_1\) there are 24 choices of families,
or 80 choices when the admissible classes are distinguished. For \(\Delta_2\) there are 24 choices in
each row of the table. For example, in the first row, with \(j=R_0,u=A_s,v=D_t,w_1=C_s,w_2=C_{-t}\), the
two outer vertices \((s+t,s),(s+t,-t)\) share row \(s+t\). In the mixed part, interchanging the copies
gives the same integrand with \(u,v\) exchanged, accounting for the factor \(2\) in the table. Finally,
\begin{align*}
 d(x)={}&G_{\rm cof}+f2E+f2z+2f2F+f3F+f3c+f4+f5c+f5F\\
        &+f6+f7+C_{3,\rm cuad}+Q_1+\Delta_1+Q_2+\Delta_2 .
\end{align*}
